% !TeX root = ../../Book4.tex
% 纲要标题：### 18.4* 无界情形
% 纲要位置：计划纲要.md，第 2971 行。

\OptionalSection{无界情形}{b4:sec:1804}

无界复形没有可供逐度归纳的起点，逐项投射或内射便不再自动足够。K-投射、K-内射与 K-平坦直接要求复形正确处理无上同调对象，并支持无界替换。

% 纲要内容（仅作写作计划，尚非教材正文）：
%
% * K-内射、K-投射与 K-平坦
% * 无界消解的必要修正
% * 不把有界证明直接无条件推广
%
% ---
%

\subsection{K-内射、K-投射与 K-平坦}
\label{b4:subsec:1804:01}

有上界投射复形和有下界内射复形之所以有效，关键不是每一项的名称，而是整个复形对正合复形的作用。无界时应直接把这种作用列为条件。

\begin{definition}\label{b4:def:k-resolutions}
设 \(R\) 为环，以下均为上链复形。
\begin{enumerate}
\item 左 \(R\)-模复形 \(P\) 若对每个正合左 \(R\)-模复形 \(A\)，都有 \(\operatorname{Hom}_R^\bullet(P,A)\) 正合，则称 \(P\) 为\term[K-toushe]{K-投射复形}[K-projective complex]。
\item 左 \(R\)-模复形 \(I\) 若对每个正合左 \(R\)-模复形 \(A\)，都有 \(\operatorname{Hom}_R^\bullet(A,I)\) 正合，则称 \(I\) 为\term[K-neishe]{K-内射复形}[K-injective complex]。
\item 左 \(R\)-模复形 \(F\) 若对每个正合右 \(R\)-模复形 \(A\)，直和总复形 \(A\otimes_RF\) 都正合，则称 \(F\) 为\term[K-pingtan]{K-平坦复形}[K-flat complex]。右模版本交换两个侧别。
\end{enumerate}
\end{definition}

前两项分别等价于所有 \(P\to A[n]\)、\(A\to I[n]\) 都零同伦，因为 Hom 复形的同调正是这些同伦类。由映射锥判据，这也等价于相应 Hom 函子把拟同构送到拟同构；K-平坦性则等价于张量函子具有这个性质。

\ProseParagraphBreak
\cref{b4:thm:bounded-projective-model}的提升论证说明，有上界的逐项投射复形是 K-投射的；其对偶说明有下界的逐项内射复形是 K-内射的。\cref{b4:lem:bounded-flat-tensor-preserves}说明有上界的逐项平坦复形是 K-平坦的。最后一个结论允许另一个变量无界，原因在于直和总化与滤过余极限的正合性，而不是对角线自动有限。

\begin{example}\label{b4:exm:unbounded-free-not-k-projective}
仍取 \(R=k[\varepsilon]/(\varepsilon^2)\)，令 \(E^n=R\) 对每个整数 \(n\)，所有微分均为乘 \(\varepsilon\)。\(E\) 正合且每项自由，但不是 K-投射：若 \(1_E\) 零同伦，则每度存在乘某元素 \(a_n\) 的同伦，使
\(1=\varepsilon a_n+a_{n+1}\varepsilon\)，右侧属于 \((\varepsilon)\)，矛盾。因此 \(\operatorname{Hom}_R(E,E)\) 的零次同调非零。

\ProseParagraphBreak
它也不是 K-平坦。取\cref{b4:exm:dual-numbers-derived-base-change}中的自由替换 \(P\to k\)。若 \(E\) 是 K-平坦，则
\(E\otimes_RP\to E\otimes_Rk\) 为拟同构。\(P\) 本身有上界且平坦，所以 \(P\otimes_RE\) 正合；交换因子的分次符号给出 \(E\otimes_RP\) 正合。但 \(E\otimes_Rk\) 每度为 \(k\)、微分为零，不正合，矛盾。

\ProseParagraphBreak
这个例子的每项甚至内射。取线性泛函 \(\lambda:R\to k\)，\(\lambda(a+b\varepsilon)=b\)。映射
\(R\to\operatorname{Hom}_k(R,k)\)、\(r\mapsto(s\mapsto\lambda(sr))\)
是 \(R\)-模同构：双线性型在基 \(1,\varepsilon\) 下的矩阵为 \(\left(\begin{smallmatrix}0&1\\1&0\end{smallmatrix}\right)\)。又有
\(\operatorname{Hom}_R(M,\operatorname{Hom}_k(R,k))\simeq\operatorname{Hom}_k(M,k)\)，右侧在向量空间短正合列上正合，故 \(R\) 内射。尽管如此，\(E\) 不是 K-内射，因为 \(E\) 正合而 \(\operatorname{Hom}_R(E,E)\) 不正合。
\end{example}

\subsection{无界消解的修正}
\label{b4:subsec:1804:02}

在模范畴中，任意无界复形仍可得到足够好的替换，但须按复形的全部循环关系构造。下面给出不依赖一般无界存在定理的方法。

\begin{definition}\label{b4:def:semifree-complex}
设 \(R\) 为环，\(P\) 是左 \(R\)-模复形。若存在子复形滤过
\[
0=P_{-1}\subseteq P_0\subseteq P_1\subseteq\cdots,
\qquad P=\bigcup_{s\ge0}P_s,
\]
使每个包含在分次模意义下分裂，每个商 \(P_s/P_{s-1}\) 是若干移位 \(R[n]\) 的直和且微分为零，则称此 \(P\) 连同滤过为\term[banziyou-fuxing]{半自由复形}[semi-free complex]。等价地，每一层可添一批齐次自由生成元，新生成元的微分只涉及更早层的生成元。
\end{definition}

\begin{theorem}\label{b4:thm:semifree-resolution}
设 \(R\) 为环，\(X\) 为任意左 \(R\)-模上链复形。存在半自由复形 \(P\) 及拟同构 \(P\to X\)。每个半自由复形都是 K-投射且 K-平坦的；这里不要求 \(X\) 或 \(P\) 有任何次数界。
\end{theorem}
\begin{proof}
先对 \(X\) 每个次数的每个余圈 \(z\in Z^nX\)，添一个次数 \(n\) 的自由生成元 \(e_z\)，令其微分为零，映到 \(z\)。所得 \(P_0\to X\) 在同调上满射。

假定已构造 \(f_s:P_s\to X\)。对每对数据
\[
z\in Z^nP_s,\qquad x\in X^{n-1},\qquad f_s(z)=\mathrm{d}_Xx,
\]
添一个次数 \(n-1\) 的自由生成元 \(e_{z,x}\)，规定
\(\mathrm{d} e_{z,x}=z\)、\(f_{s+1}(e_{z,x})=x\)。因为 \(\mathrm{d}z=0\) 且 \(f_s(z)=\mathrm{d}x\)，复形条件与映射条件同时成立。这定义 \(P_{s+1}\)，其分次模是 \(P_s\) 与所添自由模的直和，商微分为零。取并得到半自由复形 \(P\) 及映射 \(f:P\to X\)。

每个 \(X\) 的同调类在 \(P_0\) 已有余圈代表，因而 \(H(f)\) 满。若 \(z\in Z^nP\) 映成 \(X\) 中的边界，则 \(z\) 只含有限个生成元，属于某个 \(P_s\)；选定 \(x\) 使 \(f(z)=\mathrm{d}x\)，下一层就添入 \(e_{z,x}\) 使 \(z\) 成为边界。因此 \(H(f)\) 单，\(f\) 是拟同构。

为证 K-投射，设 \(A\) 正合，\(g:P\to A\) 是复形映射。逐层构造次数 \(-1\) 的映射 \(h\)，使 \(g=\mathrm{d}h+h\mathrm{d}\)。若 \(h\) 已在旧层定义，对次数 \(n\) 的新生成元 \(e\)，\(\mathrm{d}e\) 属于旧层，而且
\[
\mathrm{d}_A\bigl(g(e)-h(\mathrm{d}e)\bigr)=g(\mathrm{d}e)-\mathrm{d}_Ah(\mathrm{d}e)=0,
\]
后一等式用旧层上的同伦公式及 \(\mathrm{d}^2e=0\)。由 \(A\) 正合，可选 \(h(e)\in A^{n-1}\) 使其微分为 \(g(e)-h(\mathrm{d}e)\)。自由性允许同时线性延拓到本层。逐层取并即得全局同伦。把 \(A\) 换成任意移位，得到所有 Hom 同调消失。

为证 K-平坦，设 \(A\) 是正合右模复形。每个商
\(A\otimes_R(P_s/P_{s-1})\) 是 \(A\) 的移位直和，所以正合。滤过包含逐次分裂，张量后保持短正合列，归纳得 \(A\otimes_RP_s\) 正合。最后由滤过余极限的正合性，\(A\otimes_RP\) 正合。
\end{proof}

这里的归纳不是沿复形次数开始，而是沿添加生成元的层次开始。即使次数向两边无限延伸，每个新生成元的微分只依赖已经处理的部分，所以构造同伦仍有起点。

\subsection{有界论证向无界情形推广的条件}
\label{b4:subsec:1804:03}

\begin{proposition}\label{b4:prop:k-resolution-comparison}
设 \(R\) 为环，简记左 \(R\)-模的同伦范畴和导出范畴为 \(K(R)=K(R\text{-}\mathbf{Mod})\)、\(D(R)=D(R\text{-}\mathbf{Mod})\)。若 \(P\) 为 K-投射左 \(R\)-模上链复形，则对任意左 \(R\)-模上链复形 \(Y\)，有
\[
\operatorname{Hom}_{K(R)}(P,Y)\simeq\operatorname{Hom}_{D(R)}(P,Y).
\]
若 \(I\) 为 K-内射左 \(R\)-模上链复形，则对任意左 \(R\)-模上链复形 \(X\)，有对偶同构
\(\operatorname{Hom}_{K(R)}(X,I)\simeq\operatorname{Hom}_{D(R)}(X,I)\)。
\end{proposition}
\begin{proof}
设 \(s:Y'\to Y\) 为拟同构。\(\operatorname{Cone}(s)\) 正合，K-投射性及 Hom 的锥长正合列说明
\(\operatorname{Hom}_{K(R)}(P,Y')\to\operatorname{Hom}_{K(R)}(P,Y)\)
是同构。因此从 \(P\) 出发的同伦态射函子已经把所有分母变成同构。

更具体地，导出态射的屋顶 \(P\xleftarrow{s}U\xrightarrow{a}Y\) 中，\(s\) 对 \(\operatorname{Hom}_{K(R)}(P,-)\) 给出同构，故唯一有 \(t:P\to U\) 使 \(st=1_P\) 在同伦范畴中成立。屋顶由 \(at\) 表示。若两个这样的态射在导出范畴中相等，分式等价关系给出一个共同分母；对该分母同样施加上述 Hom 同构，便得它们已在同伦范畴中相等。对偶证明使用反向屋顶及 \(\operatorname{Hom}_{K(R)}(-,I)\)。
\end{proof}

结合\cref{b4:thm:semifree-resolution}，任意环的模范畴上可以在无界范围定义导出张量积：对两个变量取半自由替换，再用直和总化。K-平坦性使只替换一边也能计算，K-投射比较保证态射和不同替换之间的相容性。导出 Hom 则可在第一变量用半自由替换 \(P\to X\)，定义 \(\mathbf R\operatorname{Hom}_R(X,Y)=\operatorname{Hom}_R^\bullet(P,Y)\)，其中 Hom 仍取各次数的直积。这给出任意无界模复形上的计算，而不是把有限对角线的论证直接照搬。

\ProseParagraphBreak
对一般交换范畴，未必有自由生成元，也未必有所需的无限直和或精确余极限。若要用 K-内射复形构造无界 \(RF\)，必须另证每个对象有相应替换，并检查 \(F\) 在这些替换间保持拟同构；仅有“足够内射对象”并不自动提供这里使用的整个复形性质。本节不建立一般交换范畴的无界内射替换存在定理，也不让后文主干依赖它。可比较 Weibel\cite[Generalized Existence Theorem 10.5.9]{Weibel1994HomologicalAlgebra}所列出的替换与非循环性条件；有界性有限制时，前面各节已经给出了可直接使用的构造。
